% !TeX program = pdflatex
%
% "The companion by the nome: Apery's zeta(3) numbers and the modular route
%  to second solutions."  Expository.
%
% Every displayed expansion in this paper was computed exactly (Fraction
% arithmetic) by papers_out/expository_apery/verify.py; the check labels
% CHECK1..CHECK11 in section 8 are that script's output.
%
\documentclass[11pt,reqno]{amsart}

\usepackage[T1]{fontenc}
\usepackage[utf8]{inputenc}
\usepackage{amsmath,amssymb,amsthm}
\usepackage{mathtools}
\usepackage{booktabs}
\usepackage{array}
\usepackage{enumitem}
\usepackage[margin=1.05in]{geometry}
\usepackage{microtype}
\usepackage[hidelinks]{hyperref}

\theoremstyle{plain}
\newtheorem{theorem}{Theorem}[section]
\newtheorem{proposition}[theorem]{Proposition}
\newtheorem{lemma}[theorem]{Lemma}
\newtheorem{corollary}[theorem]{Corollary}

\theoremstyle{definition}
\newtheorem{definition}[theorem]{Definition}

\theoremstyle{remark}
\newtheorem{remark}[theorem]{Remark}
\newtheorem{observation}[theorem]{Observation}

\newcommand{\Z}{\mathbb{Z}}
\newcommand{\Q}{\mathbb{Q}}
\newcommand{\C}{\mathbb{C}}
\newcommand{\HH}{\mathbb{H}}
\newcommand{\Ver}[1]{\textup{\textsf{[verified: #1]}}}
\newcommand{\ta}{\theta_t}
\newcommand{\tq}{\theta_q}
\DeclareMathOperator{\ord}{ord}

\numberwithin{equation}{section}

\begin{document}

\title[The companion by the nome]
      {The companion by the nome:\\
       Ap\'ery's $\zeta(3)$ numbers and the\\ modular route to second solutions}

\author{}
\date{\today}

\subjclass[2020]{11F11, 11J72, 11B37, 33C20, 14J32}
\keywords{Ap\'ery numbers, second solution, mirror map, nome, hauptmodul,
eta quotient, Eichler integral, maximal unipotent monodromy}

\begin{abstract}
Ap\'ery's proof that $\zeta(3)\notin\Q$ runs on a three-term recurrence with
two solutions: the integers $a_n$, and a companion $b_n$ with
$b_n/a_n\to\zeta(3)$. The companion is classically described by a harmonic-sum
formula. This paper explains a different route to it --- the \emph{modular}
route --- in which the companion is read off from the nome of the associated
differential operator rather than from any sum over binomial cells. We work the
Ap\'ery case in full, because there everything is classical and checkable: the
mirror map $t(q)$ is Beukers' level-$6$ eta quotient, the operator factors as
$L=(P_3/\sigma^3)\,F\,\tq^3\,F^{-1}$, and the second solution becomes a single
Eichler-type integral,
$b_n=6\,[t^n]\,F(q)\,\tq^{-3}\bigl(t\sigma^3/(P_3F)\bigr)$.
All expansions are exact and were verified by machine over $\Q$; the ranges are
stated. A closing section explains why this detour is worth learning: in the
surrounding programme it produces companions for sporadic families where
harmonic-sum formulas provably do not exist.
\end{abstract}

\maketitle

\section{Two solutions and one miracle}

In 1978 Ap\'ery announced that
\[
  \zeta(3)=\sum_{n\ge1}\frac{1}{n^3}
\]
is irrational, and the proof --- as relayed by van der Poorten's celebrated
report \cite{vdP} --- was so compressed that the audience largely did not
believe it. The engine is one recurrence:
\begin{equation}\label{eq:rec}
  (n+1)^3u_{n+1}=(34n^3+51n^2+27n+5)\,u_n-n^3u_{n-1},
  \qquad n\ge1 .
\end{equation}
It is convenient to notice that the middle polynomial factors,
$34n^3+51n^2+27n+5=(2n+1)(17n^2+17n+5)$; the pair $(17,5)$ and the trailing
$n^3$ are the only data.

Equation~\eqref{eq:rec} has a two-dimensional solution space, and the whole
proof is the interplay of two particular members of it.

The first is
\begin{equation}\label{eq:an}
  a_n=\sum_{k=0}^{n}\binom{n}{k}^{2}\binom{n+k}{k}^{2},
  \qquad a_0=1,\ a_1=5,\ a_2=73,\ a_3=1445,\dots
\end{equation}
That these integers satisfy \eqref{eq:rec} is the first surprise (they are a
sum over $n+1$ cells, and the recurrence knows nothing about cells); it is
nowadays a two-line consequence of Zeilberger's algorithm, and we take it as
known. \Ver{$n\le 25$, exactly}

The second solution is the companion $b_n$, pinned down by
\begin{equation}\label{eq:bn}
  b_n=\sum_{k=0}^{n}\binom{n}{k}^{2}\binom{n+k}{k}^{2}
      \Biggl(\sum_{m=1}^{n}\frac{1}{m^3}
      +\sum_{m=1}^{k}\frac{(-1)^{m-1}}
        {2m^3\binom{n}{m}\binom{n+m}{m}}\Biggr),
\end{equation}
so that $b_0=0$ and $b_1=6$. We use this normalisation throughout; it is the
one in \cite{vdP}. \Ver{\eqref{eq:bn} satisfies \eqref{eq:rec} for
$1\le n\le 24$, exactly}

The miracle is that these two sequences approach each other's quotient
extremely fast. The two roots of $t^2-34t+1$ are $\alpha^{\pm1}$ with
\begin{equation}\label{eq:limit}
  \alpha=(1+\sqrt2)^4=17+12\sqrt2=33.97\ldots,
\end{equation}
these are the reciprocal singularities of \eqref{eq:rec}, and the classical
asymptotics (see \cite{vdP,Beukers79}) are
\begin{equation}\label{eq:limit2}
  a_n\sim C\,\alpha^{n}n^{-3/2},
  \qquad
  \bigl|b_n-a_n\zeta(3)\bigr|=O\bigl(\alpha^{-n}\bigr),
  \qquad\text{so}\qquad
  \lim_{n\to\infty}\frac{b_n}{a_n}=\zeta(3).
\end{equation}
Let $d_n=\mathrm{lcm}(1,\dots,n)$, so $d_n=e^{n(1+o(1))}$ by the prime number
theorem. Three inputs finish the proof.
\begin{enumerate}[label=\textup{(\roman*)},leftmargin=2.4em]
\item $a_n\in\Z$ (clear from \eqref{eq:an}) and $d_n^{3}b_n\in\Z$. The second
  is \emph{not} clear from \eqref{eq:bn} --- the inner sum there carries
  binomial coefficients in its denominators --- and is a genuine arithmetic
  theorem; see \cite[\S3]{vdP} and \cite{Beukers79}.
\item $b_n-a_n\zeta(3)\ne0$ for all large $n$. This too needs an argument; the
  cleanest is Beukers' integral representation \cite{Beukers79}, in which
  $b_n-a_n\zeta(3)$ is (up to sign) a triple integral of a strictly positive
  integrand, hence nonzero.
\item $e^{3}<\alpha$, i.e.\ $20.08\ldots<33.97\ldots$
\end{enumerate}
Now suppose $\zeta(3)=p/q$ with $p,q\in\Z$, $q\ge1$. Then
$q\,d_n^{3}\bigl(b_n-a_n\zeta(3)\bigr)$ is an integer by (i); it is nonzero by
(ii); and by \eqref{eq:limit2} and (iii) its absolute value is
$O\bigl(q\,(e^{3}/\alpha)^{n(1+o(1))}\bigr)\to0$. A nonzero integer cannot tend
to $0$. Hence $\zeta(3)\notin\Q$; that is the whole proof.

\begin{table}[t]
\centering
\begin{tabular}{@{}rrl@{}}
\toprule
$n$ & $a_n$ & $b_n$\\
\midrule
0 & $1$ & $0$\\
1 & $5$ & $6$\\
2 & $73$ & $351/4$\\
3 & $1445$ & $62531/36$\\
4 & $33001$ & $11424695/288$\\
5 & $819005$ & $35441662103/36000$\\
6 & $21460825$ & $20637706271/800$\\
\bottomrule
\end{tabular}
\qquad
\begin{tabular}{@{}rl@{}}
\toprule
$n$ & $b_n/a_n$\\
\midrule
1 & $1.2$\\
2 & $1.20205479452$\\
3 & $1.20205690119$\\
4 & $1.20205690316$\\
5 & $1.20205690316$\\
6 & $1.20205690316$\\
$\infty$ & $\zeta(3)=1.20205690316$\\
\bottomrule
\end{tabular}
\medskip
\caption{The two solutions, and the convergence.  Both columns are exact
rationals; the quotients are shown to $12$ digits. \Ver{computed}}
\label{tab:ab}
\end{table}

Table~\ref{tab:ab} is the whole story in numbers, and the reader who has never
seen it should stop and admire the right-hand column: six terms of a recurrence
give ten digits of $\zeta(3)$.

\medskip
\noindent\textbf{The question this paper answers.} Formula~\eqref{eq:bn} is a
\emph{harmonic-sum} formula: same binomial cells as $a_n$, weighted by
harmonic-type numbers. It is a beautiful thing, and for a long time it was the
only known way to write a companion down. But it is a formula of a very
particular shape, and that shape is not always available: in the surrounding
research programme (\S\ref{sec:why}) there are Ap\'ery-like families for which
one can \emph{prove} that no formula of type \eqref{eq:bn} exists, and yet the
companion is a perfectly concrete sequence with a perfectly concrete
arithmetic. We want a construction of the second solution that does not go
through cells at all.

There is one, and it is modular. The route is:
\begin{center}
recurrence $\rightsquigarrow$ differential operator $\rightsquigarrow$
Frobenius basis $\rightsquigarrow$ nome $q$ $\rightsquigarrow$ modular
parametrisation $\rightsquigarrow$ companion.
\end{center}
Every arrow is classical for Ap\'ery; only the last one is recent, and it is
the one we will derive in \S\ref{sec:punchline}. The Ap\'ery case is our worked
example precisely because the answer is already known: it lets us check every
step.

\section{From recurrence to operator}\label{sec:op}

Set $\ta=t\frac{d}{dt}$, so that $\ta t^n=nt^n$. Any three-term recurrence
whose coefficients are polynomial in $n$ becomes a differential operator by the
substitution $n\mapsto\ta$, $\text{shift}\mapsto\text{multiplication by }t$.
For \eqref{eq:rec}, with $(a,b,c,d)=(17,5,1,0)$:

\begin{definition}\label{def:L}
\[
  L \;=\; \ta^{3}\;-\;t\,(2\ta+1)(a\ta^{2}+a\ta+b)
        \;+\;t^{2}(\ta+1)\bigl(c(\ta+1)^{2}+d\bigr).
\]
\end{definition}

Expanding $L$ on a formal series $y=\sum_n u_nt^n$ and reading the coefficient
of $t^n$ gives exactly
\[
  [t^n]\,L(y)=n^3u_n-(2n-1)\bigl(a(n-1)^2+a(n-1)+b\bigr)u_{n-1}
              +(n-1)^3u_{n-2},
\]
which is \eqref{eq:rec} with $n$ replaced by $n-1$. So:

\begin{observation}\label{obs:dict}
$L(y)=0$ if and only if the coefficients of $y$ satisfy \eqref{eq:rec}
\emph{for every} $n\ge0$, including $n=0$. The $n=0$ instance reads
$1^3u_1=5u_0$, and it is the one piece of bookkeeping that will later earn its
keep: see \S\ref{sec:punchline}.
\end{observation}

Thus $y_0(t):=\sum_{n\ge0}a_nt^n$ satisfies $L(y_0)=0$
(\Ver{to $t^{26}$, exactly}), while $y_b:=\sum_n b_nt^n$ does \emph{not}:
$b$ satisfies \eqref{eq:rec} only for $n\ge1$, and the $n=0$ defect is
$1^3b_1-5b_0=6$. Hence
\begin{equation}\label{eq:inhom}
  L(y_b)=6t .
\end{equation}
This is not a blemish. It is the entire mechanism by which the second solution
is \emph{selected}: $y_0$ is the kernel, $y_b$ is the solution of an
inhomogeneous problem whose right-hand side is one monomial.

\medskip
The indicial polynomial of $L$ at $t=0$ is the $\ta$-part with $t=0$, namely
$\ta^3$: a triple root at $0$. In the language of Picard--Fuchs equations, $L$
has \emph{maximal unipotent monodromy} (MUM) at the origin. This is the
structural fact everything below rests on. It says the local solution space is
spanned by
\begin{equation}\label{eq:frob}
  y_0,\qquad
  y_1=y_0\log t+g,\qquad
  y_2=\tfrac12y_0\log^2t+g\log t+h,
\end{equation}
with $g,h\in t\Q[[t]]$ --- one holomorphic solution and two logarithmic ones,
with no fractional exponents anywhere.

Computing $g$ is elementary. Write $L=\sum_{j=0}^{2}t^jP_j(\ta)$ as in
Definition~\ref{def:L}. Because $\ta(f\log t)=(\ta f)\log t+f$ we get
$P(\ta)(f\log t)=(P(\ta)f)\log t+P'(\ta)f$, and $t^j$ commutes with all of
this, so \eqref{eq:frob} forces
\begin{equation}\label{eq:geq}
  L(g)=-\sum_{j=0}^{2}t^jP_j'(\ta)\,y_0 .
\end{equation}
Since $P_0(\ta)=\ta^3$ is invertible on $t^n$ for $n\ge1$, \eqref{eq:geq}
determines $g$ coefficientwise with no free constant beyond $g_0=0$. One turn
of the crank gives \Ver{exact}
\begin{equation}\label{eq:g}
  g(t)=12t+210t^2+4438t^3+104825t^4+\tfrac{13276637}{5}t^5+70543291\,t^6
       +\tfrac{67890874657}{35}t^7+\cdots
\end{equation}

\begin{remark}
$g$ is not $b$ and is not close to it: $g_1=12$ while $b_1=6$. The relation
between $g$ and the companion is exactly the boundary bookkeeping of
\eqref{eq:inhom}, and getting it right is the content of
\S\ref{sec:punchline}. It is worth flagging early because the naive guess
``$b_n=g_n+\lambda a_n$'' is false, and was tried.
\end{remark}

\section{The nome, and why one would invent it}

Under MUM there is a canonical coordinate. Consider the quotient $y_1/y_0$; by
\eqref{eq:frob} it is $\log t+g/y_0$, a well-defined element of
$\log t+t\Q[[t]]$. Exponentiate:

\begin{definition}[the nome]\label{def:q}
\[
  q(t)\;:=\;\exp\bigl(y_1/y_0\bigr)\;=\;t\exp\bigl(g(t)/y_0(t)\bigr)
  \;\in\; t+t^2\Q[[t]].
\]
Its compositional inverse $t(q)\in q+q^2\Q[[t]]$ is the \emph{mirror map}.
\end{definition}

Why is this the right thing to do? Two answers, one formal and one geometric.

Formally: the monodromy of $L$ around $t=0$ is $y_1\mapsto y_1+2\pi i\,y_0$,
i.e.\ $y_1/y_0\mapsto y_1/y_0+2\pi i$, i.e.\ $q\mapsto q$. The nome is the
invariant coordinate --- the coordinate in which the local monodromy has been
divided out. Whatever transcendental object is hiding behind $L$, it is a
function of $q$, not of $t$.

Geometrically: for Picard--Fuchs equations of families of Calabi--Yau type,
$y_1/y_0$ is a period ratio, so $q=e^{2\pi i\tau}$ with $\tau$ the point in the
upper half plane classifying the fibre. Then $t(q)$ is a modular function and
$y_0(t(q))$ a modular form --- \emph{if} the family is of the sort where
$\tau$ actually lives on a modular curve. For $L$ of order $3$ this is the
symmetric-square situation and it happens genuinely often; Ap\'ery is the
prototype.

Cranking out Definition~\ref{def:q} from \eqref{eq:g} \Ver{exact, to $q^{26}$}:
\begin{align}
  q(t)&=t+12t^2+222t^3+4900t^4+119133t^5+3079152t^6+\cdots,
    \label{eq:qt}\\
  t(q)&=q-12q^2+66q^3-220q^4+495q^5-804q^6+1068q^7+\cdots
    \label{eq:tq}
\end{align}

Two things should now make the reader sit up. First, $t(q)$ has \emph{integer}
coefficients --- there is no reason on earth for that in
Definition~\ref{def:q}, which involves dividing by $y_0$, exponentiating and
reverting a series, all operations that manufacture denominators freely (and
indeed $g$ itself has a $5$ and a $35$ in its denominators, \eqref{eq:g}).
Integrality of the mirror map is the fingerprint of modularity --- with the
caveat, which we will not need but which should be stated once and loudly, that
integral mirror maps are known well outside the modular-curve setting (the
quintic threefold's is integral, and no modular curve is in sight). Integrality
is \emph{evidence that prompts an identification}; the identification is what
proves something, and it is Theorem~\ref{thm:beukers} that does the work below.
Second, the
initial coefficients $1,-12,66,-220,495$ are $\binom{12}{0},-\binom{12}{1},
\binom{12}{2},-\binom{12}{3},\binom{12}{4}$: the series is trying to tell us it
is a twelfth power of something starting $1-q$.

\section{Beukers' parametrisation}\label{sec:beukers}

It is. With the Dedekind eta function
$\eta(q)=q^{1/24}\prod_{n\ge1}(1-q^n)$ (we write $q$ rather than $\tau$ in the
argument, $q=e^{2\pi i\tau}$), define
\begin{equation}\label{eq:etaquots}
  \mathbf{t}(q)=\left(\frac{\eta(q)\,\eta(q^{6})}
                           {\eta(q^{2})\,\eta(q^{3})}\right)^{\!12},
  \qquad
  \mathbf{F}(q)=\frac{\bigl(\eta(q^{2})\,\eta(q^{3})\bigr)^{7}}
                     {\bigl(\eta(q)\,\eta(q^{6})\bigr)^{5}} .
\end{equation}
The fractional powers of $q$ in the eta prefactors work out to
$q^{12(1+6-2-3)/24}=q^{1}$ for $\mathbf t$ and $q^{(5\cdot7-7\cdot5)/24}=q^{0}$
for $\mathbf F$, so $\mathbf t\in q+q^2\Z[[q]]$ and $\mathbf F\in1+q\Z[[q]]$,
as they must be. (This prefactor bookkeeping is worth doing once by hand: a
stray $q$ in front of $\mathbf t$ is the most common way to misquote
\eqref{eq:etaquots}.)

\begin{theorem}[Beukers \cite{Beukers87}]\label{thm:beukers}
$t(q)=\mathbf t(q)$ and $F(q):=y_0(t(q))=\mathbf F(q)$. Equivalently: $t$ is a
hauptmodul for the group $\Gamma_0(6)+6$ (i.e.\ $\Gamma_0(6)$ extended by the
Fricke involution $\tau\mapsto-1/6\tau$), the mirror map is the corresponding
uniformisation, and $y_0$ pulls back to a weight-$2$ modular form on
$\Gamma_0(6)$.
\end{theorem}

We do not reprove this --- it is Beukers' theorem, and the modern statement of
the surrounding picture is Zagier's \cite{Zagier09} together with
\cite{ASZ,Golyshev} --- but we do check it, because checking it is ten lines of
code and it is the identity everything else rests on.

\begin{proposition}[machine check]\label{prop:check}
With $t(q)$, $F(q)$ computed from Definition~\ref{def:q} out of the recurrence
\eqref{eq:rec} alone, and $\mathbf t,\mathbf F$ expanded from
\eqref{eq:etaquots} by the product formula, one has
$t(q)=\mathbf t(q)$ and $F(q)=\mathbf F(q)$ coefficientwise over $\Q$.
\Ver{to $q^{26}$, exact}
\end{proposition}

The first coefficients, for the reader's own ten lines:
\begin{align}
  t(q)&=q-12q^2+66q^3-220q^4+495q^5-804q^6+\cdots,\label{eq:tq2}\\
  F(q)&=1+5q+13q^2+23q^3+29q^4+30q^5+31q^6+40q^7+\cdots\label{eq:Fq}
\end{align}
Note $F$'s coefficients are small --- a weight-$2$ form has coefficients
$O(n^{1+\epsilon})$ --- whereas $a_n$ grows like $\alpha^n$. All of the growth
of the Ap\'ery numbers has been absorbed into the change of variable. That is
the practical reason this coordinate is worth the trouble: in $q$, the problem
is bounded.

\begin{remark}[the Ap\'ery numbers, recovered]
Reading Theorem~\ref{thm:beukers} backwards, $\sum a_nt^n=\mathbf F(q)$ where
$q$ is the inverse of $\mathbf t$: the Ap\'ery numbers are the Taylor
coefficients of a weight-$2$ eta quotient in the coordinate given by a weight-$0$
one. Their integrality --- obvious from \eqref{eq:an}, mysterious from
\eqref{eq:rec} --- becomes, from this side, the statement that both eta
quotients in \eqref{eq:etaquots} are integral $q$-series. Every ``$\!$Ap\'ery-like''
integrality phenomenon in the literature has, conjecturally or provably, this
shape.
\end{remark}

\section{The operator, factored through the nome}\label{sec:factor}

Now we exploit the coordinate. Write $\tq=q\frac{d}{dq}$ and
\begin{equation}\label{eq:sigma}
  \sigma\;:=\;\tq\log t\;=\;\frac{q\,t'(q)}{t(q)}\;\in\;1+q\Z[[q]],
  \qquad\text{so}\qquad \tq=\sigma\,\ta
\end{equation}
as operators on functions of $t$ (chain rule). Concretely \Ver{exact, $q^{26}$}
\begin{equation}\label{eq:sigmaexp}
  \sigma=1-12q-12q^2-12q^3-12q^4-72q^5-12q^6-96q^7-\cdots
\end{equation}
and, pleasantly, taking $\tq\log$ of \eqref{eq:etaquots} and using
$\tq\log\eta(q^m)=\tfrac{m}{24}E_2(q^m)$, where
$E_2=1-24\sum_{n\ge1}\bigl(\sum_{d\mid n}d\bigr)q^n$
is the quasimodular Eisenstein series of weight $2$, gives the closed form
\begin{equation}\label{eq:sigmaeis}
  \sigma=\tfrac12\bigl(E_2(q)+6E_2(q^{6})-2E_2(q^{2})-3E_2(q^{3})\bigr),
\end{equation}
a weight-$2$ Eisenstein series on $\Gamma_0(6)$. \Ver{to $q^{25}$, exact}
So $\sigma$ is modular data too, of the same weight as $F$.

Let $P_3(t)$ be the coefficient of $\ta^3$ in $L$, i.e.\ the leading symbol:
from Definition~\ref{def:L},
\begin{equation}\label{eq:P3}
  P_3(t)=1-2at+ct^2=1-34t+t^2,
\end{equation}
whose roots $17\pm12\sqrt2$ are (inverses of) the singular points of $L$, and
whose larger root is precisely the $\alpha$ governing the convergence rate in
\eqref{eq:limit}. Everything in Ap\'ery's proof is in this quadratic.

\subsection*{An interlude that cannot be skipped: the symmetric square}

It is tempting --- and it is wrong --- to think that the factorisation below
follows from MUM alone. Let us be exact about where the extra hypothesis
enters.

MUM gives the Frobenius shape \eqref{eq:frob}, in which $h$ is uniquely
determined; it does \emph{not} say what $h$ is. The factorisation we are after
requires the specific value
\begin{equation}\label{eq:sqcond}
  h=\frac{g^{2}}{2y_0},
  \qquad\text{equivalently}\qquad
  y_2=\frac{y_1^{2}}{2y_0},
  \qquad\text{equivalently}\qquad
  2y_0h-g^{2}=0 .
\end{equation}
For a general order-$3$ MUM operator \eqref{eq:sqcond} is false. The condition
that makes it true is that $L$ be a \emph{symmetric square}: if $u,v$ span the
solution space of a second-order operator $D$, then $u^2,uv,v^2$ span a
three-dimensional space, and the order-$3$ operator annihilating it is
$\mathrm{Sym}^2(D)$. On such an operator, $y_0=u^2$, $y_1=uv$, $y_2=v^2/2$, and
\eqref{eq:sqcond} is the identity $2\cdot u^2\cdot\tfrac{v^2}{2}=(uv)^2$ ---
trivially true, and a real restriction.

Ap\'ery's operator does satisfy it, and we prove it rather than cite it.
Normalise $L$ to be monic in $\partial=d/dt$:
\[
  \frac{1}{t^3P_3(t)}\,L=\partial^3+A_2\partial^2+A_1\partial+A_0 .
\]
The classical formula for the symmetric square of $D=\partial^2+p\partial+r$ is
\begin{equation}\label{eq:sym2}
  \mathrm{Sym}^2(D)=\partial^3+3p\,\partial^2+(2p^2+p'+4r)\,\partial
                    +(4pr+2r') .
\end{equation}
So $L$ is a symmetric square iff, setting $p:=A_2/3$ and
$r:=\tfrac14(A_1-2p^2-p')$ --- forced by matching the $\partial^2$ and
$\partial$ coefficients --- the remaining equation $A_0=4pr+2r'$ holds
identically in $\Q(t)$.

\begin{lemma}[$L$ is a symmetric square]\label{lem:sym2}
With $P_3=1-34t+t^2$, put
\[
  p=\frac{2t^2-51t+1}{t\,P_3(t)},\qquad
  r=\frac{t-10}{4t\,P_3(t)},\qquad D=\partial^2+p\,\partial+r .
\]
Then $L=t^3P_3(t)\cdot\mathrm{Sym}^2(D)$, exactly, as operators over $\Q(t)$.
Equivalently, in $\ta$-form,
\begin{equation}\label{eq:Dtheta}
  \widetilde D\;=\;\ta^{2}-t\bigl(34\ta^{2}+17\ta+\tfrac52\bigr)
                   +t^{2}\bigl(\ta+\tfrac12\bigr)^{2},
\end{equation}
which is again MUM (indicial polynomial $\ta^2$).
\end{lemma}

\begin{proof}
Expand $L$ in $\partial$: from Definition~\ref{def:L},
\[
  L=t^3(t^2-34t+1)\partial^3+3t^2(2t^2-51t+1)\partial^2
    +t(7t^2-112t+1)\partial+t(t-5).
\]
Divide by $t^3P_3$, read off $p=A_2/3$ and $r=\tfrac14(A_1-2p^2-p')$ --- these
are the values in the statement --- and verify $A_0-(4pr+2r')=0$ by clearing
denominators. This is a finite rational-function identity; it was checked
symbolically over $\Q(t)$ and the residual is identically $0$.
\Ver{exact, symbolic over $\Q(t)$; check \textup{(C12a)}}
\end{proof}

\begin{corollary}\label{cor:sq}
$y_2=y_1^2/(2y_0)$, i.e.\ \eqref{eq:sqcond} holds.
\end{corollary}

\begin{proof}
$\widetilde D$ is MUM of order $2$, so its local solutions at $0$ are
$u\in1+t\Q[[t]]$ and $v=u\log t+w$ with $w\in t\Q[[t]]$, both unique. By
Lemma~\ref{lem:sym2}, $u^2$ is a holomorphic solution of $L$ with value $1$ at
$t=0$; under MUM the holomorphic solution space of $L$ is one-dimensional, so
$u^2=y_0$. Then $uv=u^2\log t+uw=y_0\log t+uw$ is a solution of $L$ of the
Frobenius shape \eqref{eq:frob}, so by uniqueness $y_1=uv$ and $g=uw$;
and $v^2/2=\tfrac12y_0\log^2t+g\log t+\tfrac12w^2$ likewise gives $y_2=v^2/2$
and $h=w^2/2$. Hence $2y_0h-g^2=u^2w^2-(uw)^2=0$.
\end{proof}

\noindent For the reader who wants to see them,
$u=1+\tfrac52t+\tfrac{267}{8}t^2+\tfrac{10225}{16}t^3+\cdots$, and
$u^2=y_0$, $uw=g$ hold coefficientwise. \Ver{to $t^{26}$, exact; checks
\textup{(C12b--c)}}

\subsection*{The factorisation}

\begin{proposition}[the factorisation]\label{prop:factor}
Let $L$ be an order-$3$ MUM operator satisfying the symmetric-square condition
\eqref{eq:sqcond} --- for Ap\'ery's $L$ this is Corollary~\ref{cor:sq}. Then,
as operators on functions of $t$ near $t=0$,
\begin{equation}\label{eq:factor}
  L\;=\;\frac{P_3(t)}{\sigma^{3}}\;\cdot\;F\;\cdot\;\tq^{3}\;\cdot\;F^{-1},
\end{equation}
where $P_3$ is the leading symbol of $L$ in $\ta$, and $F=F(q)$, $\sigma$, $q$
are as above; $F^{\pm1}$ denote multiplication operators.
\end{proposition}

\begin{proof}
Both sides are linear differential operators of order $3$ in $t$; we show they
have the same kernel and the same leading symbol.

\emph{Kernel.} $\tq^3$ kills exactly $1,\log q,\tfrac12\log^2q$. Hence
$M:=F\tq^3F^{-1}$ kills exactly $F,\,F\log q,\,\tfrac12F\log^2q$. Now
$\log q=y_1/y_0$ by Definition~\ref{def:q}, and $F=y_0$ as functions
(Theorem~\ref{thm:beukers} is not needed here: $F$ is \emph{defined} as
$y_0\circ t$). So
\[
  F=y_0,\quad F\log q=y_0\cdot\frac{y_1}{y_0}=y_1,\quad
  \tfrac12F\log^2q=\frac{y_1^2}{2y_0}=y_2,
\]
the last step being exactly the symmetric-square condition \eqref{eq:sqcond}.
So $\ker M=\ker L$, both being the span of $y_0,y_1,y_2$. (Equivalently, and
more transparently: with $u,v$ as in Corollary~\ref{cor:sq}, $\log q=v/u$ and
$F=u^2$, so $F\cdot\{1,\log q,\tfrac12\log^2q\}=\{u^2,uv,\tfrac12v^2\}$, which
is $\ker\mathrm{Sym}^2(D)=\ker L$.)

\emph{Leading symbol.} Two order-$3$ operators with the same kernel differ by a
left multiplier, so $L=\lambda(t)\,M$ for some function $\lambda$. By
\eqref{eq:sigma}, $\tq^3=(\sigma\ta)^3=\sigma^3\ta^3+(\text{order}\le2)$, and
conjugating by the multiplication operator $F$ does not change the leading
symbol. Hence the $\ta^3$-coefficient of $M$ is $\sigma^3$, while that of $L$
is $P_3(t)$. Therefore $\lambda=P_3/\sigma^3$.
\end{proof}

The content of \eqref{eq:factor} is worth saying in words. \emph{In the nome
coordinate, and after dividing out the modular form $F$, Ap\'ery's operator is
just $\tq^3$.} All the arithmetic --- the $34n^3+51n^2+27n+5$, the $\sqrt2$,
the growth rate --- has been pushed into the change of variable $t\leftrightarrow
q$ and into the prefactor. What remains is the simplest operator of order $3$
there is.

\begin{remark}[the Yukawa coupling, and what happens without \eqref{eq:sqcond}]
\label{rem:yukawa}
It is worth knowing what \eqref{eq:factor} looks like when the symmetric-square
hypothesis fails, because the failure is not catastrophic --- it is
\emph{measured}. For a general order-$3$ MUM operator, passing to the canonical
coordinate still normalises the operator, but the normal form is
$F\,\tq\,K^{-1}\,\tq\,\tq\,F^{-1}$ (up to a left multiplier), where the extra
function $K$ is the \emph{Yukawa coupling} --- in mirror-symmetry language, the
generating series of instanton numbers, and in ours simply the obstruction
$2y_0h-g^2$ repackaged. Condition \eqref{eq:sqcond} is exactly $K$ constant, in
which case $K$ can be absorbed and the bare $\tq^3$ survives. So
Theorem~\ref{thm:companion} below is a statement about the symmetric-square
class, and the general statement carries one more datum.

This is not idle. The programme surrounding this paper finds precisely that
dichotomy empirically: all fifteen sporadic Ap\'ery-like families of
\cite{Zagier09,ASZ,Cooper} satisfy the bare formula, while the order-$5$
$\zeta(5)$ block studied in \cite{RepoZ5} does not --- it is $\mathrm{Sym}^4$-like
rather than $\mathrm{Sym}^2$-like, bare $\ta$-powers are excluded there, and it
carries a nontrivial integral Yukawa-type invariant. I thank the referee for
insisting on this distinction, which the first draft elided.
\end{remark}

And $\tq^3$ can be inverted on $q\Q[[q]]$ by inspection:
\begin{equation}\label{eq:thetainv}
  \tq^{-3}\Bigl(\sum_{m\ge1}\Psi_mq^m\Bigr)=\sum_{m\ge1}\frac{\Psi_m}{m^3}q^m .
\end{equation}
The reader will recognise the right-hand side as an \emph{Eichler integral}: a
formal $q$-series obtained from a weight-$k$ object by dividing the $m$-th
coefficient by $m^{k-1}$. Here $k-1=3$. We come back to this in \S\ref{sec:why};
for now, note that $\zeta(3)=\sum m^{-3}$ has just walked into the room by
itself.

\section{The companion, derived}\label{sec:punchline}

We can now do the thing the harmonic formula \eqref{eq:bn} does, without cells.

Recall Observation~\ref{obs:dict} and \eqref{eq:inhom}. It is cleaner to use
the \emph{unit-normalised} second solution
\[
  B(0)=0,\qquad B(1)=1,\qquad B\text{ satisfies }\eqref{eq:rec}\ (n\ge1),
\]
so that $b_n=6B(n)$ \Ver{$n\le25$, exact} and, by the same one-line defect
computation, $y_B:=\sum_nB(n)t^n$ satisfies
\begin{equation}\label{eq:LB}
  L(y_B)=t .
\end{equation}
\Ver{$L(y_B)=t$ exactly, all coefficients to $t^{26}$}

\begin{theorem}[the companion by the nome]\label{thm:companion}
Let $L$ be an order-$3$ MUM operator satisfying the symmetric-square condition
\eqref{eq:sqcond}, so that Proposition~\ref{prop:factor} applies; for Ap\'ery's
$L$ this is Lemma~\ref{lem:sym2} and Corollary~\ref{cor:sq}. Let
$F,\sigma,P_3,t(q)$ be as above and set
\begin{equation}\label{eq:Psi}
  \Psi(q)\;:=\;\frac{t(q)\,\sigma(q)^{3}}{P_3\bigl(t(q)\bigr)\,F(q)}
  \;\in\;q\Q[[q]],
  \qquad
  \Theta(q):=\tq^{-3}\Psi=\sum_{m\ge1}\frac{\Psi_m}{m^{3}}\,q^{m}.
\end{equation}
Then the second solution of \textup{\eqref{eq:rec}} is
\begin{equation}\label{eq:companion}
  B(n)=[t^{n}]\;F(q)\,\Theta(q)
  \qquad\bigl(\text{expanded in }t\text{ via }q=q(t)\bigr),
  \qquad b_n=6B(n).
\end{equation}
\end{theorem}

\begin{proof}
Substitute \eqref{eq:factor} into \eqref{eq:LB}:
$\frac{P_3}{\sigma^3}F\,\tq^3(y_B/F)=t$, i.e.\
$\tq^3(y_B/F)=t\sigma^3/(P_3F)=\Psi$. The right-hand side lies in $q\Q[[q]]$
(it has a simple zero at $q=0$ since $t=q+O(q^2)$ and
$\sigma,P_3^{-1},F^{-1}\in1+q\Q[[q]]$), so $\tq^{-3}$ applies by
\eqref{eq:thetainv} and returns the unique solution in $q\Q[[q]]$; the three
constants of integration are exactly the kernel $\{F,F\log q,\tfrac12F\log^2q\}$
of $\tq^3$ after multiplying back by $F$, and they are fixed by
$y_B\in t\Q[[t]]$ (no logarithms, no constant term). Hence $y_B=F\cdot\tq^{-3}\Psi$.
\end{proof}

That is the whole construction. It is worth restating what has been used: the
recurrence (to build $L$), MUM (to get the Frobenius basis and hence $q$), the
factorisation (Proposition~\ref{prop:factor}), and the boundary defect
\eqref{eq:LB} --- \emph{the $n=0$ instance of the recurrence, which the second
solution is precisely the one to violate}. Nothing about binomial coefficients;
nothing about harmonic numbers. If the family is modular, this runs.

\subsection*{The numbers}

Here is the computation in full, so that a reader can reproduce it.
From \eqref{eq:tq2}, \eqref{eq:Fq}, \eqref{eq:sigmaexp} and $P_3=1-34t+t^2$:
\begin{align}
  \Psi(q)&=q-19q^2+91q^3-179q^4+126q^5-q^6+344q^7-1459q^8+\cdots,
    \label{eq:Psiexp}\\
  \Theta(q)&=q-\tfrac{19}{8}q^2+\tfrac{91}{27}q^3-\tfrac{179}{64}q^4
             +\tfrac{126}{125}q^5-\tfrac{1}{216}q^6+\cdots
    \label{eq:Thetaexp}
\end{align}
(the second line is the first divided by $m^3$, term by term --- that is all
$\tq^{-3}$ is). Multiplying by $F$ and re-expanding in $t$ via \eqref{eq:qt}:

\begin{table}[ht]
\centering
\begin{tabular}{@{}rlll@{}}
\toprule
$n$ & $[t^n]F\Theta$ & $B(n)$ from \eqref{eq:rec} & $b_n=6B(n)$\\
\midrule
0 & $0$ & $0$ & $0$\\
1 & $1$ & $1$ & $6$\\
2 & $117/8$ & $117/8$ & $351/4$\\
3 & $62531/216$ & $62531/216$ & $62531/36$\\
4 & $11424695/1728$ & $11424695/1728$ & $11424695/288$\\
5 & $35441662103/216000$ & $35441662103/216000$ & $35441662103/36000$\\
6 & $20637706271/4800$ & $20637706271/4800$ & $20637706271/800$\\
\bottomrule
\end{tabular}
\medskip
\caption{Theorem~\ref{thm:companion} against the recurrence.  The two middle
columns agree for all $n\le25$, exactly over $\Q$. \Ver{$n\le25$}}
\label{tab:companion}
\end{table}

\noindent The last column of Table~\ref{tab:companion} is the last column of
Table~\ref{tab:ab}: the modular route reproduces van der Poorten's $b_n$ on the
nose. \Ver{$n\le25$, exact}

\begin{remark}[reading the answer]
The three ingredients of $\Psi$ each have a modular meaning: $t$ is the
hauptmodul, $\sigma$ is the weight-$2$ Eisenstein series \eqref{eq:sigmaeis},
$F$ is the weight-$2$ form of Theorem~\ref{thm:beukers}, and
$P_3(t)=1-34t+t^2$ is a polynomial in the hauptmodul, i.e.\ again a modular
function. Counting weights ($3\times2$ from $\sigma^3$, $-2$ from $F^{-1}$, $0$
from $t$ and $P_3$), $\Psi$ transforms with weight $4$ on $\Gamma_0(6)$.

Two honest caveats. First, $\sigma$ is quasimodular, not modular
(\eqref{eq:sigmaeis} is a combination of $E_2$'s, and only the particular
combination there --- whose $E_2$-anomalies cancel --- is genuinely weight $2$);
we have checked that cancellation only through \eqref{eq:sigmaeis}'s coefficient
identity. Second, and more to the point, ``weight $4$'' here means
\emph{meromorphic} of weight $4$: dividing by $F$ and by $P_3(t)$ can and does
introduce poles wherever those vanish, and establishing holomorphy would require
checking the zeros of $F$ and the zeros of $P_3\circ t$ against the cusps of
$\Gamma_0(6)$, which we have not done. Everything in
Theorem~\ref{thm:companion} is a statement about formal $q$-series near $q=0$,
where all the inverted quantities are units, and it is valid as such
unconditionally. The global modular reading is a (correct-looking, unverified
here) interpretation on top.

With that said, Theorem~\ref{thm:companion} reads: \emph{the companion of the
Ap\'ery numbers is $F$ times the Eichler integral of an explicit weight-$4$
object.} The little coefficients of \eqref{eq:Psiexp} are
the letters this alphabet writes with. (Amusing: $\Psi_6=-1$.)
\end{remark}

\begin{remark}[why the harmonic formula and this one are both right]
Two formulas for the same sequence must be an identity. Comparing
\eqref{eq:bn} and \eqref{eq:companion} coefficientwise gives, for each $n$, a
nontrivial identity between a cell sum with harmonic weights and a
$q$-expansion coefficient of a modular Eichler integral. The reader is invited
to write out $n=2$ ($351/4$ both ways) by hand; it is a satisfying half page,
and it explains why nobody guessed one formula from the other.
\end{remark}

\section{Why one would want the second route}\label{sec:why}

If the harmonic formula \eqref{eq:bn} already exists, the modular derivation is
a curiosity for Ap\'ery. The reason to learn it is that it survives where
\eqref{eq:bn} does not.

Ap\'ery's recurrence is one member of a now-classical zoo. Zagier's search
\cite{Zagier09} over the two-parameter family of order-$2$ analogues turned up
exactly six ``sporadic'' integral solutions; Almkvist--Zudilin
\cite{ASZ} and Cooper \cite{Cooper} extended the census at the $\zeta(3)$ level
to fifteen sporadic pairs in total. Every one of them has a second solution,
defined by the same normalisation $B(0)=0$, $B(1)=1$, and every one of them has
an Ap\'ery-type limit or an arithmetic reason not to. The natural first question
is: what is $B(n)$?

For most of the fifteen, a formula of harmonic type --- same cells, weighted by
$H$-sums --- can be found, and the working papers of the surrounding programme
\cite{RepoSporadics} give the current census of which. But for a definite
subfamily the harmonic ansatz can be \emph{closed out}: one writes down the
complete tame ansatz --- all products of harmonic sums of the right weight over
the given cells --- and proves, by exact linear algebra over enough $n$, that
the solution space is empty. For those families no formula of shape
\eqref{eq:bn} exists at all, and this is a theorem about the ansatz, not a
failure of ingenuity.

Precisely there, the construction of \S\ref{sec:punchline} works. The
programme's modular-anchor probe \cite{RepoAnchors} runs the pipeline of
\S\S\ref{sec:op}--\ref{sec:factor} on each unreachable family: it finds
that the mirror map $t(q)$ is again integral (so the family is modular), it
identifies $F(q)$ where it can --- for one such family
$F(q)=\bigl(9E_2(q^9)-E_2(q)\bigr)/8$, a weight-$2$ Eisenstein series on
$\Gamma_0(9)$; for another, $F(q)=\eta(-q)^3/\eta(-q^3)$, a sign-twisted
Borwein cubic theta --- and then applies Theorem~\ref{thm:companion} verbatim.
Out comes a companion formula in modular letters for a family where the
harmonic letters provably do not suffice. The Ap\'ery case of this paper is
that programme's control experiment: same code, known answer,
Table~\ref{tab:companion}. Related anchors and the variational route to the same
data are in \cite{RepoVar}.

The structural moral is worth a paragraph on its own. Why should a second
solution of a recurrence be an \emph{Eichler integral}? Because that is what
$\tq^{-3}$ is, and $\tq^{-3}$ is what a MUM operator becomes in its own
canonical coordinate. Eichler integrals are the classical home for objects that
are ``one integration away from modular'': $\sum a_mm^{1-k}q^m$ for $f=\sum
a_mq^m$ of weight $k$; they are not modular themselves, but their failure to be
modular is a period polynomial, and the periods are exactly the special values
one is chasing. In our case $\zeta(3)=\sum m^{-3}$ appears as the value the
Eichler integral of the trivial (Eisenstein) part contributes at the relevant
cusp, and the Ap\'ery limit \eqref{eq:limit2} should be the statement that the
$q\to$ cusp asymptotics of $F\Theta$ against $F$ is that period.
Golyshev--Zagier's ``gamma conjecture'' framework \cite{Golyshev} and the theory
of periods of the associated motive are the correct general setting; the
elementary version, which is all this paper needs, is Theorem~\ref{thm:companion}.

Let us be scrupulous about what has and has not been shown, since the previous
paragraph is the one place where enthusiasm could outrun the argument.
\emph{This paper constructs the companion modularly; it does not derive the
limit $b_n/a_n\to\zeta(3)$ from modular transformation theory.} Doing that would
mean evaluating the Eichler integral $\Theta$ at the relevant cusp and
identifying the resulting period with $\zeta(3)$ --- a genuine computation in
the theory of periods of Eichler integrals, and one we have not performed. The
limit \eqref{eq:limit2} is used here exactly as Ap\'ery used it: as a classical
fact about the recurrence. The modular story explains the \emph{shape} of the
companion; the $\zeta(3)$ itself remains, in this account, an input.

The practical upshot for a working reader: if you have an integral
Ap\'ery-like recurrence and want the second solution, do not start by guessing
harmonic weights. Compute $g$, compute $q$, look at whether $t(q)$ is integral.
If it is, you are ten lines of code from the companion, and those ten lines are
in the appendix.

\section{The checks}\label{sec:checks}

Nothing in this paper was asserted without computing it. All arithmetic is
exact over $\Q$ (Python \texttt{Fraction}s); series are truncated at $q^{26}$,
$t^{26}$. The script is \texttt{verify.py} alongside this file; its output is:

\begin{enumerate}[label=\textup{(C\arabic*)},leftmargin=3.2em]
\item $a_n$ from \eqref{eq:an} equals $a_n$ from \eqref{eq:rec}, $n\le25$. PASS
\item $b_n$ from \eqref{eq:bn} satisfies \eqref{eq:rec} for $1\le n\le24$, and
      $b_0=0$, $b_1=6$; the $n=0$ defect is $1^3b_1-5b_0=6$. PASS
\item $b_{25}/a_{25}$ agrees with $\zeta(3)$ to $40$ decimal digits. PASS
\item $t(q)=\bigl(\eta_1\eta_6/\eta_2\eta_3\bigr)^{12}$ to $q^{26}$. PASS
\item $F(q)=(\eta_2\eta_3)^7/(\eta_1\eta_6)^5$ to $q^{26}$. PASS
\item $B(n)=[t^n]F\,\tq^{-3}\Psi$ equals the recurrence's $B$, $n\le26$;
      stated in the text as $n\le25$ because $q^{26}$ is the truncation edge of
      the $\sigma$ construction. PASS
\item $b_n=6B(n)$, $n\le25$. PASS
\item $L(y_B)=t$ exactly: coefficient $1$ at $t^1$, zero at every other
      $t^n$, $n\le26$. PASS
\item $L(y_0)=0$ to $t^{26}$. PASS
\item the eta prefactor bookkeeping of \eqref{eq:etaquots} (exponents
      $q^1$ and $q^0$). PASS
\item $\sigma=\tfrac12(E_2(q)+6E_2(q^6)-2E_2(q^2)-3E_2(q^3))$ to $q^{25}$. PASS
\item[\textup{(C12a)}] the symmetric-square residual $A_0-(4pr+2r')$ of
      Lemma~\ref{lem:sym2} is identically $0$ in $\Q(t)$ --- symbolic, not
      truncated. PASS
\item[\textup{(C12b)}] $u^2=y_0$ to $t^{26}$, where $u$ is the holomorphic
      solution of $\widetilde D$. PASS
\item[\textup{(C12c)}] $u\,w=g$ to $t^{26}$, where $w$ is $\widetilde D$'s
      log-tail; hence $h=w^2/2$ and $2y_0h-g^2=0$. PASS
\end{enumerate}

Everything labelled \textup{\textsf{[verified: \dots]}} in the text refers to
this list. Theorem~\ref{thm:beukers} is quoted from the literature; (C4)--(C5)
are a check of it, not a proof. Propositions~\ref{prop:factor} and
Theorem~\ref{thm:companion} are proved in the text, and (C6)--(C8) are their
independent numerical confirmation. (C12a) is the only check that is a proof
rather than evidence: it is an identity between rational functions, verified
symbolically, and Lemma~\ref{lem:sym2} rests on it. (C12b--c) are consistency
checks on Corollary~\ref{cor:sq}, which is itself proved.

\appendix

\section{Ten lines of code}

The following is the skeleton; \texttt{smul}, \texttt{sinv}, \texttt{sexp},
\texttt{srevert}, \texttt{compose} are the obvious truncated power-series
operations on lists of \texttt{Fraction}s (product, reciprocal, exponential,
compositional inverse, composition), each five lines.

\begin{small}
\begin{verbatim}
from fractions import Fraction as F
N = 26; a,b,c,d = 17,5,1,0

# 1. the Apery numbers, from the recurrence
A = [F(1), F(5)]
for n in range(1, N+2):
    A.append((F((2*n+1)*(a*n*n+a*n+b))*A[n] - F(n*(c*n*n+d))*A[n-1])
             / F((n+1)**3))
y0 = A[:N+1]

# 2. the log-solution tail g:  L(g) = -sum_j t^j P_j'(theta) y0
#    with P_0 = th^3, P_1 = -(2th+1)(a th^2 + a th + b),
#         P_2 = (th+1)(c(th+1)^2 + d)          [solve coefficientwise]
g = gseries(y0)

# 3. the nome and the mirror map
qt = smul([F(0),F(1)]+[F(0)]*(N-1), sexp(smul(g, sinv(y0))))  # q(t)
tq = srevert(qt)                                              # t(q)
Fq = compose(y0, tq)                                          # F(q)

# 4. sigma = theta_q log t,  P3 = 1 - 2a t + c t^2
T   = [tq[i+1] for i in range(N)] + [F(0)]                    # t/q
sig = smul([F(i)*T[i] for i in range(N+1)], sinv(T)); sig[0] += F(1)
t2  = smul(tq, tq)
P3  = [F(-2*a)*tq[i] + F(c)*t2[i] for i in range(N+1)];  P3[0] = F(1)

# 5. the companion
Psi   = smul(smul(tq, smul(sig, smul(sig, sig))),
             smul(sinv(P3), sinv(Fq)))
Theta = [F(0)] + [Psi[m]/F(m)**3 for m in range(1, N+1)]
Bn    = compose(smul(Fq, Theta), qt)     # Bn[n] = B(n);  b_n = 6*B(n)
\end{verbatim}
\end{small}

\noindent Step 2 is the only one needing a word: writing $L=\sum_jt^jP_j(\ta)$,
one solves \eqref{eq:geq} for $g$ by reading off $t^N$, using that $P_0(N)=N^3
\ne0$ for $N\ge1$:
\[
  g_N=\frac{1}{N^3}\Bigl(R_N-\sum_{j\ge1}P_j(N-j)\,g_{N-j}\Bigr),
  \qquad
  R=-\sum_j t^jP_j'(\ta)y_0 .
\]
Substituting steps 1--5 for any other Ap\'ery-like $(a,b,c,d)$ is the whole
generalisation.

\begin{thebibliography}{99}
\sloppy

\bibitem{ASZ} G.~Almkvist, W.~Zudilin,
\emph{Differential equations, mirror maps and zeta values},
in: Mirror Symmetry V, AMS/IP Stud.\ Adv.\ Math.\ \textbf{38} (2006), 481--515.

\bibitem{Beukers87} F.~Beukers,
\emph{Irrationality proofs using modular forms},
Ast\'erisque \textbf{147--148} (1987), 271--283.

\bibitem{Beukers79} F.~Beukers,
\emph{A note on the irrationality of $\zeta(2)$ and $\zeta(3)$},
Bull.\ London Math.\ Soc.\ \textbf{11} (1979), 268--272.

\bibitem{Cooper} S.~Cooper,
\emph{Sporadic sequences, modular forms and new series for $1/\pi$},
Ramanujan J.\ \textbf{29} (2012), 163--183.

\bibitem{Golyshev} V.~Golyshev, D.~Zagier,
\emph{Proof of the gamma conjecture for Fano 3-folds of Picard rank one},
Izv.\ Math.\ \textbf{80} (2016), 24--49.

\bibitem{vdP} A.~van der Poorten,
\emph{A proof that Euler missed\dots{} Ap\'ery's proof of the irrationality of
$\zeta(3)$}, Math.\ Intelligencer \textbf{1} (1979), 195--203.

\bibitem{Zagier09} D.~Zagier,
\emph{Integral solutions of Ap\'ery-like recurrence equations},
in: Groups and Symmetries, CRM Proc.\ Lecture Notes \textbf{47} (2009), 349--366.

\bibitem{RepoAnchors} \emph{Modular anchors for the deformation-unreachable
sporadic families}, working paper of the present programme
(\texttt{papers\_out/modular\_anchors}; probe notes
\texttt{work/MODULAR\_COACTION\_PROBE.md}, scripts
\texttt{work/z5eps/eps48\_modular\_nome.py},
\texttt{eps49\_zeta\_companion.py}).

\bibitem{RepoZ5} \emph{Modularity and Yukawa integrality of the $\zeta(5)$
five-block}, working notes of the present programme
(\texttt{work/Z5\_MODULARITY\_PROBE.md}).

\bibitem{RepoVar} \emph{The variational anchor}, working paper of the present
programme (\texttt{papers\_out/variational\_anchor}).

\bibitem{RepoSporadics} \emph{The fifteen sporadic Ap\'ery-like pairs: second
solutions, the $\chi$-twisted Lucas law, and the closed forms}, working paper of
the present programme (\texttt{papers\_out/sporadics}).

\end{thebibliography}

\end{document}
